% Pass5404--5409: all-q unsigned apartment visibility spectrum.
\subsection{Unsigned apartment visibility: the full flag space before Hodge signing}

Let \(B=|C|\) be the unsigned flag-edge/apartment incidence matrix from the preceding
apartment theorem.  The distance-\(d\) pair count is \(q^{4-d}\), so
\[
\boxed{
BB^{\mathsf T}=q^4A_0+q^3A_1+q^2A_2+qA_3+A_4.
}
\]
This is the unsigned companion to
\[
CC^{\mathsf T}=q^4A_0-q^3A_1+q^2A_2-qA_3+A_4.
\]
The sign change is not cosmetic: it changes the rank from the protected cycle sector to the
entire flag space.

Evaluating the unsigned radial kernel on the five adjacency eigenspaces of the flag distance
scheme gives
\[
\boxed{
\begin{aligned}
\lambda_0&=8q^4,\\
\lambda_+&=2q^2(q-1)\bigl(q+1+\sqrt{2q}\bigr),\\
\lambda_m&=2q^2(q-1)^2,\\
\lambda_-&=2q^2(q-1)\bigl(q+1-\sqrt{2q}\bigr),\\
\lambda_c&=(q-1)^2(q^2+1),
\end{aligned}}
\]
with multiplicities
\[
\boxed{
1,\quad \frac{q(q+1)^2}{2},\quad q(q^2+1),\quad
\frac{q(q+1)^2}{2},\quad q^4.
}
\]
For \(q>1\), every eigenvalue is positive.  In particular,
\[
\boxed{\operatorname{rank}B=N=(q+1)^2(q^2+1).}
\]
The complete unsigned apartment family therefore spans the full real flag space.

By contrast, the signed matrix satisfies
\[
\operatorname{rank}C=q^4.
\]
The missing dimensions are exactly
\[
\boxed{
N-q^4
=2(q+1)(q^2+1)-1
=\operatorname{rank}D,
}
\]
where \(D\) is the oriented incidence matrix of the connected Levi graph.  Thus the
alternating point-to-line cycle signs remove precisely the Kirchhoff cut-space sector and
retain precisely the cycle-space sector:
\[
\boxed{
\text{unsigned visibility}=\text{cycle}\oplus\text{cut},
\qquad
\text{signed visibility}=\text{cycle}.
}
\]
This supplies an all-\(q\) explanation of the q=3 rank complement \(160=81+79\).

At \(q=3\), the unsigned spectrum becomes
\[
\boxed{
648^1
\oplus(144+36\sqrt6)^{24}
\oplus72^{30}
\oplus(144-36\sqrt6)^{24}
\oplus40^{81},
}
\]
exactly the BT546 unsigned Levi edge/apartment spectrum.  The signed matrix has rank \(81\),
while the removed cut sector has rank \(79\), exactly as in BT549.  Hence the old 3-adic
visibility kernel and the old cycle/cut tight-frame decomposition are both specializations of
the same all-\(q\) orientation-sign mechanism.

\paragraph{Evidence boundary.}
This is a combinatorial spectral theorem about the full unsigned apartment-incidence Gram
matrix and its relation to the signed Hodge matrix.  It is not a contextuality-degree theorem,
a quantum-code minimum-distance theorem, or a physical noise-tolerance result.
